\section{三角函数线}

\begin{definition}[任意角的三角函数定义]
    设角 $\alpha$ 的顶点为坐标原点，始边为 $x$ 轴的非负半轴，终边上任意一点 $P$ 的坐标为 $(x, y)$，它与原点的距离为 $r = \sqrt{x^2+y^2}$ ($r>0$)，则：
    $$\sin\alpha = \frac{y}{r}, \quad \cos\alpha = \frac{x}{r}, \quad \tan\alpha = \frac{y}{x}$$
    \textbf{全是天才：一句话记忆象限}
\end{definition}

\begin{figure}[H]
    \centering
    \begin{tikzpicture}[scale=1]
        \draw (0,0) circle (2);
        \coordinate (origo) at (0,0);
        \coordinate[label = above left:$P$] (P) at (1,{sqrt(3)});
        \draw[thick,->] (origo) -- ++(3,0) node[black,right] (xaxis){$x$};
        \draw[thick,gray,->] (origo) -- ++(0,3) node (mary) [black,above] {$y$};
        \draw[thick] (origo) -- ++(60:2) coordinate (bob);
        \pic [draw, ->, red, angle eccentricity=2] {angle = xaxis--origo--bob} node[left, color=red] {$\theta$};
        \draw[blue, ->] (1,{sqrt(3)}) -- (1,0);
        \draw[cyan, ->] (1,{sqrt(3)}) -- (0,{sqrt(3)});
        \node[color=cyan] at (-0.5,{sqrt(3)}) {$\cos \theta$};
        \node[color=blue] at (1,-0.3) {$\sin \theta$};
        \draw[black,dotted] (-2,2)--(3,2);
        \draw[black,dotted] ({2*sqrt(3)/3},2)--(1,{sqrt(3)});
        \draw[brown,->] ({2*sqrt(3)/3},2)--({2*sqrt(3)/3},0);
        \node[color=brown] at (1.6,1) {$\tan \theta$};
        \node at (P)[circle,fill,inner sep=2pt]{};
    \end{tikzpicture}
    \caption{三角函数线}
\end{figure}

\subsection{反三角函数}

\begin{definition}[反三角表示法]
    指的是当某个角的特殊值不能被直接算出来时，使用 $arc$ 来表示角。
\end{definition}

\begin{problem}
\begin{itemize}
    \item 请问 $arccos\frac{1}{2}$ 表示的角是什么？
    \item 请问 $arctanln2$ 表示什么？
\end{itemize}
\end{problem}

\begin{example}
    已知角 $\alpha$ 的终边经过点 $P(-3, 4)$，求 $\sin\alpha$, $\cos\alpha$, $\tan\alpha$ 的值。
\end{example}

\begin{solution}
    点 $P(-3, 4)$ 中，$x=-3, y=4$，所以 $r = \sqrt{(-3)^2 + 4^2} = \sqrt{9+16} = \sqrt{25} = 5$。
    $$\sin\alpha = \frac{y}{r} = \frac{4}{5}, \quad \cos\alpha = \frac{x}{r} = -\frac{3}{5}, \quad \tan\alpha = \frac{y}{x} = -\frac{4}{3}$$
\end{solution}

\vspace{2em}
